% 第 2 题图甲：手抛撒谷粒的简笔示意（TikZ 矢量重绘，替代原人物插画）
% 人物仅作情景示意，题目所需的轨迹信息见图乙
\begin{tikzpicture}[font=\small, line width=0.9pt, line cap=round, scale=0.85]
  % 水田
  \draw (-1.7,0.28) .. controls (-0.6,0.5) and (0.7,0.48) .. (1.9,0.3);
  \draw (-1.2,0.05) -- (-0.95,0.05);
  \draw (-0.5,0.02) -- (-0.25,0.02);
  \draw (0.35,0.03) -- (0.6,0.03);
  \draw (1.0,0.08) -- (1.25,0.08);
  % 人体
  \draw (0.05,1.86) -- (0,1.15);
  \draw (0,1.15) -- (-0.24,0.36);
  \draw (0,1.15) -- (0.24,0.4);
  % 手臂
  \draw (0.03,1.72) -- (-0.26,1.36);
  \draw (0.03,1.72) -- (0.44,1.56);
  % 头与斗笠
  \fill (0.05,2.02) circle (0.13);
  \draw (-0.3,2.14) -- (0.4,2.14) -- (0.05,2.4) -- cycle;
  \draw (-0.18,2.14) -- (-0.18,2.2);
  % 簸箕
  \draw (0.4,1.68) -- (0.74,1.74) -- (0.77,1.5) -- (0.43,1.44) -- cycle;
  % 谷粒
  \foreach \p in {(0.86,1.62),(0.95,1.52),(1.04,1.42),(0.88,1.42),
                  (1.12,1.3),(0.98,1.28),(1.2,1.14),(1.06,1.1),
                  (1.28,0.98),(1.12,0.94),(1.36,0.82),(1.2,0.78),
                  (1.44,0.66),(1.28,0.62),(1.5,0.5),(1.36,0.46)} {
    \fill \p circle (0.028);}
  \node[below] at (0.05,-0.12) {甲};
\end{tikzpicture}
